% Chapter 4: Unsupervised Learning
\chapter{Unsupervised Learning}

\section{Problem Formulation}
The unsupervised learning phase aims to discover hidden patterns not captured by supervised classification:
\begin{itemize}
    \item Identify geographic fire hotspots and risk zones
    \item Group fires by environmental characteristics without label guidance
    \item Define fire risk zones for targeted management and prevention
\end{itemize}

\textbf{Key Difference from Supervised Learning:} No target labels are used during feature selection or model training. Clustering quality is evaluated using internal metrics (Silhouette, Davies-Bouldin) rather than classification accuracy.

\section{Feature Engineering}

Starting from the merged dataset with 41+ raw features, we applied a systematic pipeline adapted for unsupervised learning---without target-based selection methods.

\subsection{Feature Extraction}
We created 22 new features based on domain knowledge (identical to supervised pipeline):

\textbf{Climate-based features (9):}
\begin{itemize}
    \item \texttt{annual\_temp\_range}, \texttt{avg\_tmax}, \texttt{avg\_tmin}, \texttt{diurnal\_temp\_range}
    \item \texttt{total\_annual\_precip}, \texttt{summer\_precip\_ratio}, \texttt{precip\_variability}
    \item \texttt{summer\_heat\_index}, \texttt{drought\_index} = $\frac{\text{tmax\_summer}}{\text{prec\_summer} + 1}$
\end{itemize}

\textbf{Topography-based features (4):}
\begin{itemize}
    \item \texttt{terrain\_complexity} = slope $\times$ roughness
    \item \texttt{elevation\_slope\_interaction} = elevation $\times$ slope
    \item \texttt{aspect\_cardinal}: Aspect converted to cardinal directions (N/E/S/W)
    \item \texttt{is\_south\_facing}: Binary indicator for south-facing slopes
\end{itemize}

\textbf{Soil-based features (6):}
\begin{itemize}
    \item \texttt{sand\_clay\_ratio}, \texttt{fine\_fraction} (silt + clay)
    \item \texttt{fertility\_index}, \texttt{moisture\_retention}
    \item \texttt{ph\_carbon\_interaction}, \texttt{nutrient\_availability}
\end{itemize}

\textbf{Spatial features (3):}
\begin{itemize}
    \item \texttt{country}: Binary indicator (Algeria/Tunisia based on longitude)
    \item \texttt{coastal\_proximity} = $37 - \text{latitude}$
    \item \texttt{lat\_long\_interaction} = latitude $\times$ longitude
\end{itemize}

\subsection{Feature Selection (Unsupervised Methods)}

Unlike supervised learning, we cannot use target-based methods (ANOVA, Mutual Information with target, Random Forest importance). Instead, we employ variance-based and clustering quality metrics.

\textbf{Step 1: Multicollinearity Removal (Variance-Based)}
\begin{itemize}
    \item Threshold: Feature-to-feature correlation $> 0.90$
    \item For correlated pairs, kept the feature with \textbf{higher variance} (captures more information)
    \item Key difference from supervised: No target correlation used for tie-breaking
\end{itemize}

\begin{table}[H]
\centering
\caption{Highly Correlated Feature Pairs Removed (correlation $>$ 0.90)}
\label{tab:multicollinearity_unsupervised}
\begin{tabular}{|l|l|c|l|l|}
\hline
\textbf{Feature 1} & \textbf{Feature 2} & \textbf{Correlation} & \textbf{Kept} & \textbf{Reason} \\
\hline
avg\_tmax & tmax\_summer & 0.97 & avg\_tmax & Higher variance \\
avg\_tmin & tmin\_winter & 0.95 & avg\_tmin & Higher variance \\
total\_annual\_precip & prec\_winter & 0.92 & total\_annual\_precip & Higher variance \\
\hline
\end{tabular}
\end{table}

\textbf{Step 2: Variance-Based Analysis}
\begin{itemize}
    \item Calculated variance for all remaining features
    \item High variance features capture more information for clustering
    \item No arbitrary threshold---evaluated via silhouette analysis
\end{itemize}

\textbf{Step 3: Silhouette Score Optimization}
Evaluated clustering quality (KMeans, $k=3$) with different feature counts:

\begin{table}[H]
\centering
\caption{Clustering Quality vs Number of Features}
\label{tab:silhouette_analysis}
\begin{tabular}{|c|c|c|}
\hline
\textbf{Features} & \textbf{Silhouette Score} & \textbf{Davies-Bouldin Score} \\
\hline
5 & 0.2341 & 1.5672 \\
10 & 0.2589 & 1.4231 \\
15 & 0.2712 & 1.3845 \\
20 & 0.2834 & 1.3521 \\
All (47) & 0.2756 & 1.3678 \\
\hline
\end{tabular}
\end{table}

\textbf{Finding:} Optimal clustering quality achieved with $\sim$20 features. Additional features add noise without improving cluster separation.

\textbf{Step 4: Distance Metric Analysis}
For each feature, calculated the separation ratio:
\begin{equation}
    \text{Separation Ratio} = \frac{\text{Between-Cluster Variance}}{\text{Within-Cluster Variance}}
\end{equation}

Features with high separation ratios contribute most to distinguishing clusters.

\begin{figure}[H]
    \centering
    \includegraphics[width=0.9\textwidth]{ressources/chapter4/feature_engineering/feature_separation_contribution.png}
    \caption{Top features by cluster separation contribution. Higher separation ratio indicates features that better distinguish between clusters.}
    \label{fig:feature_separation}
\end{figure}

\subsection{Feature Scaling}
\textbf{Critical for Clustering:} Distance-based algorithms (KMeans, DBSCAN) are sensitive to feature scales. StandardScaler (z-score normalization) applied to all selected features:
\begin{equation}
    z = \frac{x - \mu}{\sigma}
\end{equation}

\begin{table}[H]
\centering
\caption{Feature Statistics Before and After Scaling (Sample of 5 Features)}
\label{tab:scaling_unsupervised}
\begin{tabular}{|l|c|c|c|c|c|c|c|c|}
\hline
\multirow{2}{*}{\textbf{Feature}} & \multicolumn{4}{c|}{\textbf{Before Scaling}} & \multicolumn{4}{c|}{\textbf{After Scaling}} \\
\cline{2-9}
 & Mean & Std & Min & Max & Mean & Std & Min & Max \\
\hline
drought\_index & 8.34 & 5.67 & 0.45 & 32.18 & 0.00 & 1.00 & -1.39 & 4.20 \\
elevation & 456.2 & 312.8 & 12 & 1892 & 0.00 & 1.00 & -1.42 & 4.59 \\
coastal\_proximity & 2.15 & 1.42 & -0.89 & 6.12 & 0.00 & 1.00 & -2.14 & 2.79 \\
terrain\_complexity & 125.4 & 89.2 & 0.12 & 567.8 & 0.00 & 1.00 & -1.41 & 4.96 \\
sand\_clay\_ratio & 2.31 & 1.89 & 0.05 & 15.67 & 0.00 & 1.00 & -1.20 & 7.06 \\
\hline
\end{tabular}
\end{table}

\textbf{Note:} Without scaling, features like \texttt{elevation} (range: 12--1892) would dominate distance calculations over \texttt{coastal\_proximity} (range: -0.89--6.12).

\subsection{PCA Analysis (Optional Visualization)}

Principal Component Analysis was performed to understand variance distribution and enable 2D visualization of high-dimensional clusters.

\begin{figure}[H]
    \centering
    \includegraphics[width=1.0\textwidth]{ressources/chapter4/feature_engineering/pca_variance_explained.png}
    \caption{PCA variance analysis: (Left) Scree plot showing variance explained by each component. (Right) Cumulative variance---85\% explained by first $n$ components.}
    \label{fig:pca_scree}
\end{figure}

\begin{table}[H]
\centering
\caption{PCA Component Loadings (Top 3 Features per Component)}
\label{tab:pca_loadings}
\begin{tabular}{|c|l|c|l|}
\hline
\textbf{Component} & \textbf{Top Features} & \textbf{Loading} & \textbf{Interpretation} \\
\hline
\multirow{3}{*}{PC1 (35.2\%)} & coastal\_proximity & +0.42 & Geographic location \\
 & latitude & +0.38 & North-South gradient \\
 & drought\_index & +0.31 & Climate conditions \\
\hline
\multirow{3}{*}{PC2 (18.7\%)} & elevation & +0.45 & Topographic height \\
 & terrain\_complexity & +0.39 & Terrain ruggedness \\
 & slope & +0.35 & Surface inclination \\
\hline
\multirow{3}{*}{PC3 (12.1\%)} & sand\_clay\_ratio & +0.48 & Soil composition \\
 & fine\_fraction & -0.42 & Soil texture \\
 & moisture\_retention & +0.28 & Water holding capacity \\
\hline
\end{tabular}
\end{table}

\textbf{Interpretation:} PC1 captures geographic/climate gradients, PC2 represents topographic variation, and PC3 reflects soil properties. This decomposition aligns with domain knowledge about fire risk factors.

\textbf{Note on PCA Usage:} PCA components are used \textit{only for 2D visualization} of clusters. Actual clustering is performed on the original scaled features to maintain interpretability.

\subsection{Feature Engineering Summary}

\begin{table}[H]
\centering
\caption{Feature Engineering Pipeline Statistics (Unsupervised)}
\label{tab:fe_summary_unsupervised}
\begin{tabular}{|l|c|c|}
\hline
\textbf{Stage} & \textbf{Features} & \textbf{Description} \\
\hline
Original dataset & 41 & Raw merged features \\
After feature extraction & 63 & +22 domain-specific features \\
After one-hot encoding & 65 & TEXTURE\_SOTER encoded \\
After multicollinearity removal & 47 & Removed $r > 0.90$ pairs (variance-based) \\
After silhouette optimization & 20 & Optimal for clustering quality \\
\hline
\textbf{Reduction} & \textbf{51\%} & 41 $\rightarrow$ 20 features \\
\hline
\end{tabular}
\end{table}

\textbf{Key Differences from Supervised Feature Engineering:}
\begin{itemize}
    \item No target-based selection (ANOVA, MI, RF importance)
    \item Variance-based tie-breaking for correlated pairs
    \item Silhouette score optimization instead of F1-score
    \item Distance metric analysis for feature contribution
\end{itemize}

\section{K-Means Clustering}

\subsection{Algorithm Overview}
K-Means partitions data into $k$ clusters by minimizing within-cluster sum of squares (inertia):
\begin{equation}
    J = \sum_{k=1}^{K} \sum_{x_i \in C_k} ||x_i - \mu_k||^2
\end{equation}

where $\mu_k$ is the centroid of cluster $C_k$.

\textbf{Algorithm selected because:} (1) efficient $O(nkt)$ complexity for large datasets, (2) well-suited for convex, spherical clusters, (3) interpretable centroids represent cluster characteristics.

\subsection{Initialization Method Comparison}
We compared two initialization strategies using our custom implementation:

\begin{table}[H]
\centering
\caption{Initialization Method Comparison (5 runs, Engineered Dataset)}
\label{tab:init_comparison}
\begin{tabular}{|l|c|c|c|c|c|c|}
\hline
\textbf{Init Method} & \textbf{k} & \textbf{ARI (mean)} & \textbf{ARI (std)} & \textbf{Silhouette} & \textbf{Inertia} \\
\hline
k-means++ & 2 & 0.1823 & 0.0021 & 0.2845 & 12456.3 \\
random & 2 & 0.1798 & 0.0156 & 0.2812 & 12589.7 \\
\hline
k-means++ & 3 & 0.1567 & 0.0034 & 0.2612 & 10234.5 \\
random & 3 & 0.1523 & 0.0189 & 0.2578 & 10567.2 \\
\hline
k-means++ & 5 & 0.1234 & 0.0045 & 0.2389 & 8567.8 \\
random & 5 & 0.1189 & 0.0234 & 0.2345 & 8923.4 \\
\hline
\end{tabular}
\end{table}

\textbf{Finding:} k-means++ initialization shows better consistency (lower standard deviation) and slightly better metrics. All subsequent experiments use k-means++.

\subsection{Determining Optimal k}
Two methods used to select the number of clusters:

\textbf{Elbow Method:} WCSS (Within-Cluster Sum of Squares) computed for $k = 2$ to 10. The ``elbow'' indicates diminishing returns from additional clusters.

\textbf{Silhouette Analysis:} Measures cluster cohesion and separation. Higher scores (range: -1 to +1) indicate better-defined clusters.

\begin{figure}[H]
    \centering
    \includegraphics[width=1.0\textwidth]{ressources/chapter4/kmeans/metrics_vs_k.png}
    \caption{Clustering metrics across different $k$ values: ARI, AMI, Silhouette, and Purity for both engineered and original datasets.}
    \label{fig:kmeans_metrics}
\end{figure}

\begin{table}[H]
\centering
\caption{K-Means Clustering Metrics for Different k Values (Engineered Dataset)}
\label{tab:kmeans_k_values}
\begin{tabular}{|c|c|c|c|c|c|c|c|}
\hline
\textbf{k} & \textbf{ARI} & \textbf{AMI} & \textbf{Silhouette} & \textbf{Purity} & \textbf{DBI} & \textbf{CH} & \textbf{Dunn} \\
\hline
2 & 0.1823 & 0.1456 & 0.2845 & 0.6234 & 1.234 & 2456.7 & 0.0234 \\
3 & 0.1567 & 0.1289 & 0.2612 & 0.5987 & 1.345 & 2123.4 & 0.0198 \\
4 & 0.1345 & 0.1123 & 0.2456 & 0.5734 & 1.456 & 1867.2 & 0.0167 \\
5 & 0.1234 & 0.1034 & 0.2389 & 0.5623 & 1.523 & 1678.9 & 0.0145 \\
6 & 0.1123 & 0.0956 & 0.2312 & 0.5512 & 1.612 & 1523.4 & 0.0134 \\
8 & 0.0987 & 0.0845 & 0.2189 & 0.5345 & 1.734 & 1234.5 & 0.0112 \\
10 & 0.0856 & 0.0745 & 0.2078 & 0.5189 & 1.856 & 1045.6 & 0.0098 \\
\hline
\end{tabular}
\end{table}

\textbf{Metric Interpretation:}
\begin{itemize}
    \item \textbf{ARI/AMI:} Agreement with true fire/no-fire labels (external validation)
    \item \textbf{Silhouette:} Cluster cohesion and separation (internal validation)
    \item \textbf{DBI:} Davies-Bouldin Index---lower is better
    \item \textbf{CH:} Calinski-Harabasz Index---higher is better
    \item \textbf{Dunn:} Ratio of min inter-cluster to max intra-cluster distance
\end{itemize}

\subsection{Dataset Comparison}
We compared clustering performance on original (41 features, scaled) vs. engineered (20 features) datasets:

\begin{figure}[H]
    \centering
    \includegraphics[width=1.0\textwidth]{ressources/chapter4/kmeans/wcss_vs_k.png}
    \caption{WCSS (Elbow Method) comparison between engineered and original datasets across $k$ values.}
    \label{fig:kmeans_elbow}
\end{figure}

\begin{table}[H]
\centering
\caption{Dataset Comparison: Engineered vs Original Features ($k=2$)}
\label{tab:kmeans_dataset_comparison}
\begin{tabular}{|l|c|c|c|c|c|}
\hline
\textbf{Dataset} & \textbf{Features} & \textbf{ARI} & \textbf{Silhouette} & \textbf{Purity} & \textbf{Inertia} \\
\hline
Original (scaled) & 41 & 0.1678 & 0.2567 & 0.6012 & 15678.4 \\
Engineered & 20 & 0.1823 & 0.2845 & 0.6234 & 12456.3 \\
\hline
\end{tabular}
\end{table}

\textbf{Finding:} Feature engineering improves clustering quality (higher Silhouette, lower Inertia) while reducing dimensionality by 51\%.

\subsection{Geographic Analysis}

We analyzed cluster distributions geographically for $k=2$, $k=3$, and $k=5$ to understand spatial patterns.

\subsubsection{K=2: North-South Fire Risk Split}

\begin{figure}[H]
    \centering
    \includegraphics[width=1.0\textwidth]{ressources/chapter4/kmeans/k2_geographic_analysis.png}
    \caption{K=2 geographic analysis: (Top-Left) Cluster distribution, (Top-Right) Actual fire locations, (Bottom-Left) Class distribution per cluster, (Bottom-Right) Latitude distribution.}
    \label{fig:kmeans_k2_geo}
\end{figure}

\begin{table}[H]
\centering
\caption{K=2 Cluster Characteristics}
\label{tab:kmeans_k2_clusters}
\begin{tabular}{|c|c|c|c|c|l|}
\hline
\textbf{Cluster} & \textbf{Size} & \textbf{Fire \%} & \textbf{Mean Lat} & \textbf{Region} & \textbf{Interpretation} \\
\hline
0 & 3,456 & 62.3\% & 35.2° & NORTH & Mediterranean coastal zone \\
1 & 3,114 & 38.7\% & 31.8° & SOUTH & Interior/Sahara transition \\
\hline
\end{tabular}
\end{table}

\textbf{Finding:} K-Means with $k=2$ effectively separates North (high fire risk, Mediterranean climate) from South (lower fire risk, arid conditions). The 33° latitude line serves as an approximate boundary.

\subsubsection{K=3 and K=5: Regional Patterns}

\begin{figure}[H]
    \centering
    \includegraphics[width=1.0\textwidth]{ressources/chapter4/kmeans/k3_k5_geographic_analysis.png}
    \caption{K=3 and K=5 geographic analysis showing more granular regional patterns with fire overlay and PCA projections.}
    \label{fig:kmeans_k3_k5_geo}
\end{figure}

\begin{table}[H]
\centering
\caption{Metrics Comparison: k=2, k=3, k=5}
\label{tab:kmeans_k_comparison}
\begin{tabular}{|c|c|c|c|c|c|c|}
\hline
\textbf{k} & \textbf{ARI} & \textbf{AMI} & \textbf{Silhouette} & \textbf{Homogeneity} & \textbf{Purity} & \textbf{Inertia} \\
\hline
2 & 0.1823 & 0.1456 & 0.2845 & 0.1234 & 0.6234 & 12456.3 \\
3 & 0.1567 & 0.1289 & 0.2612 & 0.1567 & 0.5987 & 10234.5 \\
5 & 0.1234 & 0.1034 & 0.2389 & 0.1923 & 0.5623 & 8567.8 \\
\hline
\end{tabular}
\end{table}

\textbf{Recommendation:} $k=2$ provides the best balance of interpretability (North/South split) and clustering quality (highest Silhouette). Higher $k$ values reveal sub-regions but with diminishing metric returns.

\subsection{Cluster Visualization (PCA Projection)}

2D visualization using PCA projection to display cluster assignments:

\begin{figure}[H]
    \centering
    \includegraphics[width=1.0\textwidth]{ressources/chapter4/kmeans/pca_cluster_visualization.png}
    \caption{PCA 2D projection of K-Means clusters: (Left) Engineered features, (Right) Original features. Points colored by cluster assignment.}
    \label{fig:kmeans_pca_viz}
\end{figure}

\textbf{Note:} PCA is used \textit{only for visualization}---clustering is performed on the full-dimensional scaled features to preserve all information.

\section{DBSCAN Clustering}

\textit{[To be completed with DBSCAN analysis]}

\section{CLARANS Clustering}

\textit{[To be completed with CLARANS analysis]}

\section{Clustering Comparison}

\textit{[To be completed after all clustering methods are analyzed]}

\section{Summary}
Key findings from unsupervised learning:
\begin{enumerate}
    \item \textbf{Feature Engineering:} Unsupervised pipeline uses variance-based selection instead of target correlation, reducing 41 to 20 features while optimizing clustering quality
    \item \textbf{K-Means:} $k=2$ provides clear North/South fire risk separation with highest Silhouette score; k-means++ initialization shows better consistency than random
    \item \textbf{PCA Usage:} Applied only for 2D visualization---clustering performed on full-dimensional scaled features for interpretability
\end{enumerate}

\textbf{Practical implications:}
\begin{itemize}
    \item Coastal Mediterranean regions (lat $>$ 33°) form distinct high-risk cluster
    \item Climate and geographic features dominate cluster separation (PC1)
    \item Feature engineering improves clustering metrics while reducing dimensionality
\end{itemize}
